\documentclass{article}
\usepackage[T1]{fontenc}
\usepackage{lmodern}
\usepackage{graphicx}
\usepackage{amsmath,amssymb}
\usepackage[a4paper,margin=2cm]{geometry}
\usepackage[absolute,overlay]{textpos}
\setlength{\TPHorizModule}{1mm}
\setlength{\TPVertModule}{1mm}
\usepackage[indonesian]{babel}
\usepackage{hyperref}
\setlength{\emergencystretch}{2em}
% unit: Halaman 8

\begin{document}

% === Halaman 8 (python/native) ===

• Secara de-fakto, hanya ada dua varian AES, yaitu AES-128 dan AES-

256, karena akan sangat jarang pengguna menggunakan kunci yang 

panjangnya 192 bit.


Panjang Kunci 

(Nk words) 

Ukuran Blok 

(Nb words) 

Jumlah Putaran 

(Nr) 

AES-128 

4 

4 

10 

AES-192 

6 

4 

12 

AES-256 

8 

4 

14 


Catatan: 1 word = 32 bit

\end{document}
